% Source PDF page 675 (printed p. 661)
本附录简要介绍数值线性代数的一些基础内容，着重讨论求解一个或多个线性方程组的方法。我们主要介绍直接法，即非迭代方法，以及如何利用问题结构提高效率。数值线性代数还有许多重要问题与方法不在这里讨论，包括数值稳定性、矩阵分解的细节、适用于并行或多处理器的方法，以及迭代方法。关于这些及其他主题，读者可参阅本附录末尾列出的文献。

\section{矩阵结构与算法复杂度}
我们着重讨论线性方程组
\begin{equation}
 Ax=b\tag{C.1}
\end{equation}
的求解方法，其中 $A\in\R^{n\times n}$、$b\in\R^n$。假设 $A$ 非奇异，因此对任意 $b$，解都唯一，且为 $x=A^{-1}b$。这一基本问题出现在许多优化算法中，而且往往占据大部分计算量。在线性方程组(C.1)的求解中，矩阵 $A$ 常称为\emph{系数矩阵}，向量 $b$ 称为\emph{右端向量}。

求解(C.1)的标准通用方法，计算量大致按 $n^3$ 增长。这些方法除要求 $A$ 非奇异外，不对它作其他假设，因此适用范围很广。当 $n$ 为几百或更小时，除要求最苛刻的实时应用外，通用方法可能是最好的选择。当 $n$ 超过一千左右时，求解 $Ax=b$ 的通用方法就不那么实用了。

% Source PDF page 676 (printed p. 662)
\subsection*{系数矩阵的结构}
许多情况下，系数矩阵 $A$ 具有某种特殊结构或形式，可以采用针对这种结构的方法，更高效地求解 $Ax=b$。例如，在牛顿方程组 $\nabla^2f(x)\Delta x_{\mathrm{nt}}=-\nabla f(x)$ 中，系数矩阵对称正定，因此可以使用比通用方法快约一倍的方法，而且其舍入误差性质也更好。还有许多其他可利用的结构，它们节省的计算量，或算法加速比，通常远不止两倍。许多情况下，计算量可由通用方法的 $n^3$ 降低至正比于 $n^2$，甚至正比于 $n$。这些方法通常用于 $n$ 至少为一百、而且常常远大于一百的情况，因此节省的计算量可能非常可观。

可以利用的系数矩阵结构多种多样。与稀疏模式，即矩阵中零元素和非零元素的分布有关的简单例子，包括带状矩阵、分块对角矩阵和稀疏矩阵。一种更不明显的可利用结构，是对角矩阵加低秩矩阵。许多常见形式的凸优化问题，都会产生系数矩阵具有这些结构的线性方程组。还有许多其他可利用的矩阵结构，例如 Toeplitz、Hankel 和循环矩阵，本附录不作讨论。

不利用矩阵中任何稀疏模式的通用方法，称为针对\emph{稠密矩阵}的方法；完全不利用任何矩阵结构的方法，称为针对\emph{无结构矩阵}的方法。

\subsection{用浮点运算次数分析复杂度}
数值线性代数算法的代价，常用执行算法所需的\emph{浮点运算}总次数表示，简称 \emph{flop 数}，并将其写成各个问题维数的函数。一次 flop 定义为对两个浮点数进行一次加、减、乘或除运算。一些作者将一次乘法及紧随其后的一次加法合称为一次 flop，因此他们的计数最多会小一倍。分析算法复杂度时，我们统计总 flop 数，将它表示为相关矩阵和向量维数的函数，通常是多项式，再只保留最高阶或占主导地位的项，略去其余各项。

例如，设某个算法总共需要
\[
 m^3+3m^2n+mn+4mn^2+5m+22
\]
次浮点运算，其中 $m,n$ 为问题维数。通常把它简化为
\[
 m^3+3m^2n+4mn^2
\]
次，因为这些是关于 $m,n$ 的最高阶项。若进一步假设 $m\ll n$，则可继续简化为 $4mn^2$。

% Source PDF page 677 (printed p. 663)
浮点运算计数最初流行时，浮点运算相对较慢，因此计数能够很好地估计总计算时间。现在的情况已经不同：缓存边界、访问局部性等因素，都可能显著影响数值算法的计算时间。不过，flop 数仍能粗略估计数值算法的计算时间，以及计算时间如何随问题规模增长。由于 flop 数已不能准确预测算法的计算时间，我们通常最关注其阶数，即最高次数，而忽略小于两倍左右的差别。例如，flop 数为 $5n^2$ 的算法，可以看作与 $4n^2$ 的算法相当，但比 $(1/3)n^3$ 的算法快。

\subsection{基本矩阵与向量运算的代价}
\subsubsection*{向量运算}
计算两个向量 $x,y\in\R^n$ 的内积 $x^Ty$，先计算各个乘积 $x_iy_i$，再求和，需要 $n$ 次乘法与 $n-1$ 次加法，即 $2n-1$ 次浮点运算。如上所述，只保留最高阶项，就说内积需要 $2n$ 次浮点运算；更粗略地说，所需浮点运算次数为 $n$ 量级。标量与向量相乘 $\alpha x$，其中 $\alpha\in\R$、$x\in\R^n$，需要 $n$ 次浮点运算。两个向量 $x,y\in\R^n$ 相加 $x+y$，同样需要 $n$ 次。

若向量 $x,y$ 稀疏，即只有少量非零元素，并使用适当的数据结构存储，就可以更快地完成这些基本运算。例如，若稀疏向量 $x$ 有 $N$ 个非零元素，则计算内积 $x^Ty$ 需要 $2N$ 次浮点运算。

\subsubsection*{矩阵与向量相乘}
计算矩阵向量乘积 $y=Ax$，其中 $A\in\R^{m\times n}$，需要 $2mn$ 次浮点运算：需要求出 $y$ 的 $m$ 个分量，每个分量都是 $A$ 的一行与 $x$ 的内积，即两个 $\R^n$ 中向量的内积。

利用 $A$ 的结构，常可加快矩阵向量乘法。例如，若 $A$ 是对角矩阵，计算 $Ax$ 只需 $n$ 次浮点运算，而一般 $n\times n$ 矩阵的乘法需要 $2n^2$ 次。更一般地，若 $A$ 稀疏，在全部 $mn$ 个元素中只有 $N$ 个非零元素，则计算 $Ax$ 需要 $2N$ 次，因为可以跳过与零有关的乘法和加法。

一个不那么明显的例子是：矩阵 $A$ 的秩为 $p\ll\min\{m,n\}$，并以分解形式 $A=UV$ 表示和存储，其中 $U\in\R^{m\times p}$、$V\in\R^{p\times n}$。这时可先计算 $Vx$，需要 $2pn$ 次，再计算 $U(Vx)$，需要 $2mp$ 次，总共为 $2p(m+n)$ 次浮点运算。由于 $p\ll\min\{m,n\}$，这个数远小于 $2mn$。

\subsubsection*{矩阵与矩阵相乘}
计算矩阵乘积 $C=AB$，其中 $A\in\R^{m\times n}$、$B\in\R^{n\times p}$，需要 $2mnp$ 次浮点运算。需要求出 $C$ 的 $mp$ 个元素，每个元素都是两个长度为 $n$ 的
% Source PDF page 678 (printed p. 664)
向量的内积。同样，利用 $A,B$ 的结构，往往可以节省大量计算。例如，若 $A,B$ 稀疏，就可以跳过与零有关的加法和乘法，从而加快计算。若 $m=p$，且已知 $C$ 对称，则只需计算下三角部分的 $(1/2)m(m+1)$ 个元素，矩阵乘法便只需 $m^2n$ 次浮点运算。

计算多个矩阵的乘积时，可以采用不同的两两相乘次序，而这些次序通常有不同的 flop 数。最简单的例子是 $D=ABC$，其中 $A\in\R^{m\times n}$、$B\in\R^{n\times p}$、$C\in\R^{p\times q}$。通过矩阵两两相乘，可以用两种方法计算 $D$。第一种方法先算 $AB$，需要 $2mnp$ 次，再算 $D=(AB)C$，需要 $2mpq$ 次，总共为 $2mp(n+q)$ 次。另一种方法先算 $BC$，需要 $2npq$ 次，再算 $D=A(BC)$，需要 $2mnq$ 次，总共为 $2nq(m+p)$ 次。当 $2mp(n+q)<2nq(m+p)$，即
\[
 \frac1n+\frac1q<\frac1m+\frac1p,
\]
第一种方法更好。这里假设矩阵相乘时没有利用任何矩阵结构。

超过三个矩阵的乘积，有许多不同的两两相乘次序。给定矩阵维数，设计一个确定最佳次序，即所需 flop 数最少的次序的算法并不难；不过，在大多数应用中，最佳次序是明确的。

\section{利用矩阵分解求解线性方程组}
\subsection{容易求解的线性方程组}
先考察几种容易求解 $Ax=b$ 的情形，也就是容易计算 $x=A^{-1}b$ 的情形。

\subsubsection*{对角矩阵}
设 $A$ 为非奇异对角矩阵，即对所有 $i$ 有 $a_{ii}\neq0$。线性方程组 $Ax=b$ 可以写成 $a_{ii}x_i=b_i$，$i=1,\ldots,n$。解为 $x_i=b_i/a_{ii}$，计算需要 $n$ 次浮点运算。

\subsubsection*{下三角矩阵}
若 $A\in\R^{n\times n}$ 满足 $j>i$ 时 $a_{ij}=0$，就称 $A$ 为\emph{下三角矩阵}。对角元素都等于一的下三角矩阵，称为\emph{单位下三角矩阵}。下三角矩阵非奇异，当且仅当对所有 $i$ 都有 $a_{ii}\neq0$。

% Source PDF page 679 (printed p. 665)
设 $A$ 为非奇异下三角矩阵，则 $Ax=b$ 为
\[
 \begin{bmatrix}
 a_{11}&0&\cdots&0\\
 a_{21}&a_{22}&\cdots&0\\
 \vdots&\vdots&\ddots&\vdots\\
 a_{n1}&a_{n2}&\cdots&a_{nn}
 \end{bmatrix}
 \begin{bmatrix}x_1\\x_2\\\vdots\\x_n\end{bmatrix}
 =\begin{bmatrix}b_1\\b_2\\\vdots\\b_n\end{bmatrix}.
\]
由第一行，$a_{11}x_1=b_1$，所以 $x_1=b_1/a_{11}$。第二行为 $a_{21}x_1+a_{22}x_2=b_2$，因此 $x_2=(b_2-a_{21}x_1)/a_{22}$。$x_1$ 已经算出，所以右端每个数都已知。继续这一过程，就可以用前面的分量表示 $x$ 的每个分量，得到算法
\begin{align*}
 x_1&\coloneqq b_1/a_{11},\\
 x_2&\coloneqq(b_2-a_{21}x_1)/a_{22},\\
 x_3&\coloneqq(b_3-a_{31}x_1-a_{32}x_2)/a_{33},\\
 &\ \vdots\\
 x_n&\coloneqq(b_n-a_{n1}x_1-a_{n2}x_2-\cdots-a_{n,n-1}x_{n-1})/a_{nn}.
\end{align*}
这个过程称为\emph{前代}：把已知值代入下一个方程，依次求出 $x$ 的分量。

计算前代的 flop 数如下：先求 $x_1$，需要一次；将 $x_1$ 代入第二个方程求 $x_2$，需要三次；再将 $x_1,x_2$ 代入第三个方程求 $x_3$，需要五次，依此类推。总次数为
\[
 1+3+5+\cdots+(2n-1)=n^2.
\]
因此，$A$ 为非奇异下三角矩阵时，计算 $x=A^{-1}b$ 需要 $n^2$ 次浮点运算。

若 $A$ 除下三角外还具有其他结构，前代所需运算可少于 $n^2$ 次。例如，若 $A$ 稀疏或为带状矩阵，且每行至多有 $k$ 个非零元素，则每一步前代至多需要 $2k+1$ 次浮点运算，总 flop 数为 $2(k+1)n$；略去 $2n$ 项后，为 $2kn$。

\subsubsection*{上三角矩阵}
若 $A^T$ 是下三角矩阵，即 $j<i$ 时 $a_{ij}=0$，则称 $A\in\R^{n\times n}$ 为\emph{上三角矩阵}。系数矩阵为非奇异上三角矩阵时，可用类似前代的方法求解，只是先求 $x_n$，再求 $x_{n-1}$，依此类推。算法为
\begin{align*}
 x_n&\coloneqq b_n/a_{nn},\\
 x_{n-1}&\coloneqq(b_{n-1}-a_{n-1,n}x_n)/a_{n-1,n-1},\\
 x_{n-2}&\coloneqq(b_{n-2}-a_{n-2,n-1}x_{n-1}-a_{n-2,n}x_n)/a_{n-2,n-2},\\
 &\ \vdots\\
 x_1&\coloneqq(b_1-a_{12}x_2-a_{13}x_3-\cdots-a_{1n}x_n)/a_{11}.
\end{align*}
% Source PDF page 680 (printed p. 666)
这称为\emph{后代}或\emph{回代}，因为各分量按逆序求出。通过回代计算 $x=A^{-1}b$ 需要 $n^2$ 次浮点运算。若 $A$ 为稀疏或带状上三角矩阵，每行至多有 $k$ 个非零元素，则回代需要 $2kn$ 次。

\subsubsection*{正交矩阵}
若 $A\in\R^{n\times n}$ 满足 $A^TA=I$，即 $A^{-1}=A^T$，就称 $A$ \emph{正交}。这时只需计算矩阵向量乘积 $x=A^Tb$，即可得到 $x=A^{-1}b$，一般需要 $2n^2$ 次浮点运算。

若 $A$ 还具有其他结构，计算 $x=A^{-1}b$ 的运算数还可少于 $2n^2$。例如，若 $A=I-2uu^T$，$\norm{u}_2=1$，则可计算
\[
 x=A^{-1}b=(I-2uu^T)^Tb=b-2(u^Tb)u.
\]
先求 $u^Tb$，再计算 $b-2(u^Tb)u$，需要 $4n$ 次浮点运算。

\subsubsection*{置换矩阵}
设 $\pi=(\pi_1,\ldots,\pi_n)$ 是 $(1,2,\ldots,n)$ 的一个置换。相应的\emph{置换矩阵} $A\in\R^{n\times n}$ 为
\[
 A_{ij}=\begin{cases}1,&j=\pi_i,\\0,&\text{其他情况}.\end{cases}
\]
置换矩阵的每一行或每一列都恰有一个元素等于一，其余元素都为零。向量乘以置换矩阵，只是把各分量重新排列：
\[
 Ax=(x_{\pi_1},\ldots,x_{\pi_n}).
\]
置换矩阵的逆，是逆置换 $\pi^{-1}$ 对应的置换矩阵，结果等于 $A^T$，因此置换矩阵是正交矩阵。

当 $A$ 为置换矩阵时，$Ax=b$ 很容易求解：按 $\pi^{-1}$ 对 $b$ 的分量作置换，就得到 $x$。按照本书对浮点运算的定义，这不需要任何浮点运算，但具体实现可能涉及浮点数的复制。从 $x=A^Tb$ 也能得到同一结论：与 $A$ 一样，$A^T$ 的每行只有一个值为一的非零元素，因此无需加法，所需乘法也都只是乘以一。

\subsection{分解—求解法}
求解 $Ax=b$ 的基本方法，是将 $A$ 表示为非奇异矩阵的乘积：
\[
 A=A_1A_2\cdots A_k,
\]
% Source PDF page 681 (printed p. 667)
于是
\[
 x=A^{-1}b=A_k^{-1}A_{k-1}^{-1}\cdots A_1^{-1}b.
\]
按此公式从右向左计算，可以得到 $x$：
\begin{align*}
 z_1&\coloneqq A_1^{-1}b,\\
 z_2&\coloneqq A_2^{-1}z_1=A_2^{-1}A_1^{-1}b,\\
 &\ \vdots\\
 z_{k-1}&\coloneqq A_{k-1}^{-1}z_{k-2}=A_{k-1}^{-1}\cdots A_1^{-1}b,\\
 x&\coloneqq A_k^{-1}z_{k-1}=A_k^{-1}\cdots A_1^{-1}b.
\end{align*}
第 $i$ 步需要计算 $z_i=A_i^{-1}z_{i-1}$，即求解线性方程组 $A_iz_i=z_{i-1}$。如果每个方程组都容易求解，例如 $A_i$ 为对角、下三角、上三角或置换矩阵，这就给出了一种计算 $x=A^{-1}b$ 的方法。

把 $A$ 写成分解形式，即计算各个因子 $A_i$，称为\emph{分解步骤}；通过依次求解 $A_iz_i=z_{i-1}$ 这一类问题，递推算出 $x=A^{-1}b$，则常称为\emph{求解步骤}。分解—求解法求解 $Ax=b$ 的总 flop 数为 $f+s$，其中 $f$ 为计算分解的 flop 数，$s$ 为求解步骤的总 flop 数。许多情况下，分解代价 $f$ 主导总的求解代价 $s$，这时求解 $Ax=b$，即计算 $x=A^{-1}b$ 的代价，就只计为 $f$。

\subsubsection*{多个右端向量的方程组}
假设需要求解
\[
 Ax_1=b_1,\qquad Ax_2=b_2,\qquad\ldots,\qquad Ax_m=b_m,
\]
其中 $A\in\R^{n\times n}$ 非奇异。也就是说，需要求解 $m$ 个具有同一系数矩阵、但右端向量不同的线性方程组。也可以把它看作计算矩阵
\[
 X=A^{-1}B,
\]
其中
\[
 X=\begin{bmatrix}x_1&x_2&\cdots&x_m\end{bmatrix}\in\R^{n\times m},\qquad
 B=\begin{bmatrix}b_1&b_2&\cdots&b_m\end{bmatrix}\in\R^{n\times m}.
\]
为此，先分解 $A$，代价为 $f$。然后对 $i=1,\ldots,m$，利用求解步骤计算 $A^{-1}b_i$。由于只分解 $A$ 一次，总计算量为
\[
 f+ms.
\]
换句话说，分解代价分摊在 $m$ 次求解上。若对每个 $i$ 都无谓地重复分解步骤，代价就会变成 $m(f+s)$。

当分解代价 $f$ 主导求解代价 $s$ 时，分解—求解法求解少量具有同一系数矩阵的线性方程组，代价基本上与求解一个方程组相同。这是因为最昂贵的分解步骤只进行一次。
